% ============================================
% Supplementary Note 7: Alternative Models and Controls for Community-Level Selection
% ============================================
\subsection*{\rev{Supplementary Note 7: Alternative models and controls for community-level selection}}
\label{sec:suppl:alternative_models_controls}

\rev{Several alternative mechanisms could in principle generate origin-correlated persistence during coalescence, including passive biomass loading, initial-richness effects, and environmental modification by dominant taxa. Because coalescence was initiated by equal-volume pooling without OD normalization, heterogeneity in parental-community biomass at mixing could in principle bias the coalesced community toward the higher-biomass parent. To test this suggested hypothesis, we asked whether parental-community OD$_{600}$ correlates with the parental dominance direction or with pairwise selection correlation. If Dominance were primarily driven by unequal absolute biomass at mixing, the higher-OD parental community should win more frequently, and signed $\Delta\mathrm{OD}_{A-B}$, the difference between the OD of parental community A and parental community B, should be positively associated with PDI. We observed the opposite pattern. The higher-OD parental community won only 37\% of Base Dominance events and 13.2\% of Nutr$+$ Dominance events, signed $\Delta\mathrm{OD}_{A-B}$ was weakly related to PDI in Base and strongly negatively associated with PDI in Nutr$+$, and OD-binned within-community pairwise selection correlation did not increase consistently across media (Supplementary Figs.~35--36). Thus, parental-community biomass differences do not explain the nutrient-dependent increase in Dominance or correlated species fates.}

\rev{We also tested whether initial species-pool size could account for the observed community-level selection signatures in the experiments and simulations (Supplementary Fig.~34). In the synthetic-community experiments, we did not detect a significant effect of initial pool size on Dominance frequency across 6, 12, and 24 species ($\chi^2 = 2.24$, $p = 0.69$), and we also did not detect a pool-size effect within individual nutrient conditions (Dominance vs non-Dominance tests: Nutr$-$ $p = 0.39$, Base $p = 0.82$, Nutr$+$ $p = 0.27$). Initial pool size affected assembly structure, with realized parental richness increasing with pool size and ASV richness per inoculated species decreasing with pool size, but these assembly-richness effects did not translate into a corresponding pool-size dependence of Dominance frequency.}

\rev{In the gLV model, we performed a pool-size ablation at $\mu \in \{0.30, 0.60, 0.80\}$, varying the initial pool size from 4 to 48 species per parental community. Dominance frequency was primarily set by $\mu$ and changed little with pool size, whereas the same-parental-community versus cross-parental-community pairwise-selection gap increased with pool size. Reconstructed post-assembly interaction matrices showed lower within-community than between-community interaction strengths across pool sizes, supporting the interpretation that pool size affects assembly filtering and correlated species fates more than it affects Dominance-class frequency. Thus, pool-size variation does not explain the main transition in community-level selection; instead, it modulates the strength of correlated species fates within the interaction-strength regimes described in the main text.}

\rev{Finally, we considered environmental-feedback models motivated by pH-mediated microbial interactions \citep{Ratzke2018,Ratzke2020}. A Ratzke--Gore-style pH-feedback model can generate asymmetric replacement over an intermediate feedback range, showing that environmental modification alone can contribute to Dominance-like outcomes (Supplementary Fig.~44). However, in our implementation this pH-only model did not reproduce the monotone high-interaction Dominance trend as directly as the effective-interaction gLV framework. Together with the empirical pH-pair analysis (Supplementary Fig.~37), this supports pH modification as a plausible contributor to winner identity in Nutr$+$, rather than as a complete explanation for the observed Nutr$+$ Dominance frequency or as a replacement for the outcome-level community-level selection interpretation.}

\rev{We also tested whether pH mismatch itself increased Dominance frequency by comparing acid--alk parental-community pairs with pooled same-pH pairs. Dominance frequency was not significantly increased in either Base or Nutr$+$ medium (Fisher's exact tests, $p = 0.49$ and $p = 0.47$), although within acid--alk pairs the acidic parental community dominated most Nutr$+$ events (38/44, 86.4\%; binomial $p = 9.4 \times 10^{-7}$; Supplementary Fig.~37).}

\rev{Two additional controls address related versions of the same concern. First, because dominant taxa could in principle create circularity in the PDI analysis, we recalculated PDI after removing either the most abundant species in the coalesced community or the dominant species from each parental community. The Nutr$+$ association weakened but remained significant, whereas Base did not retain a significant association (Supplementary Fig.~33). This supports the interpretation that dominant species contribute to the Nutr$+$ top-down-associated regime, but that the Nutr$+$ signal is not solely an artifact of including those same dominant species in the PDI calculation.}

\rev{Second, abundance-skew and richness-dependent geometric effects can create apparent Dominance under simple null models. We address this separately in Supplementary Note~1 and Supplementary Figs.~38--39 using event-matched additive null models and simulation parent-vector norm analyses. In the experiment, the event-matched additive null was usually classified as Mixture rather than Dominance, including 66/83 Base medium events, and the observed Base outcomes were shifted further toward parental asymmetry than their matched additive nulls (one-sided paired Wilcoxon test, $p = 5.9 \times 10^{-14}$; Supplementary Fig.~38). In the simulations, increasing $\mu$ reduced final richness in both post-assembly parental communities and coalesced communities, and the simple additive null consequently produced more Dominance at high $\mu$ (0.0\%, 5.7\%, and 16.2\% at $\mu = 0.30, 0.60,$ and $0.80$, respectively; Supplementary Fig.~39). However, the full gLV coalescence simulations produced much higher Dominance at the same interaction strengths (20.4\%, 61.0\%, and 69.5\%). Thus, the low-richness/norm-imbalance effect is real and points in the direction raised by the reviewer, but it is marginal in the experimental data and quantitatively insufficient to explain the simulated Dominance transition.}

\rev{Finally, the baseline gLV model does not represent facilitative interactions explicitly. The facilitative-tail, reciprocal pair-coupling, weak-mutualism, and mean-vs-variance analyses are therefore retained in Supplementary Note~3 and Supplementary Figs.~40--43, where they define the scope of the effective-interaction framework. These analyses do not imply that facilitation is absent in all coalescence systems; rather, they show that the main Dominance trend is not restricted to the narrowest iid competitive-matrix assumption.}
